% Generated ablation-bars.pdf is supplied by export_ablation_evaluation.py.
\begin{figure*}[tbh]
\centering
\includegraphics[width=\textwidth]{ablation-bars.pdf}
\caption{Ablation comparison at 3,000 source nodes: (a) valid delivery per source opportunity, (b) conditional valid-report delay, and (c) modeled invalid acceptance per admitted report. Bars show equally weighted means over 20 paired seeds; error bars denote pointwise 95\% percentile-bootstrap confidence intervals from 100,000 shared seed resamples. Source opportunities include suppressed reports. Delay is conditional on valid delivery. Modeled invalid acceptance is not a cryptographic forgery probability. All recovery records are right-censored, so no mean recovery time or recovery-time ranking is inferred.}
\label{fig:ablation_archived}
\end{figure*}

% Replacement for the previous numerical recovery-time interpretation.
Fig.~\ref{fig:ablation_archived} and Table~\ref{tab:performance_costs}
show service--cost trade-offs. The static without-Markov policy delivers
44.010\% of source opportunities versus 35.572\% for DA-AHE, while its
conditional valid-report delay is 10.840~s versus 6.859~s and its modeled
CPU cost is 1.367~J/run versus 0.822~J/run. Their admission fractions are
100\% and 52.603\%, respectively, so this comparison concerns complete
policies. Without IntDb2 yields 33.395\% valid delivery per opportunity
and a conditional delay of 3.625~s. Without route repair coincides with
DA-AHE in the archive. Disabling HVE results in 2.087\% modeled invalid
acceptance among admitted reports while valid service remains unchanged.
All recovery outcomes are right-censored, precluding a recovery-time ranking.
